\documentclass[a4paper,twocolumn,pdflatex,ja=standard]{bxjsarticle}
% a4jでは画像挿入時に問題が起こる場合があるため，a4paperを使用する．

% ===== 基本 =====
\geometry{
    reset,
    a4paper,
    left=2.0cm,
    right=2.0cm,
    top=2.0cm,
    bottom=2.0cm
}
\usepackage{amsmath,amssymb,mathtools,bm}
\usepackage{graphicx}
\usepackage{xcolor}
\usepackage{url}
\usepackage{wrapfig}
\usepackage{enumitem}
\usepackage{listings}
\usepackage{booktabs}
\usepackage{cuted}
\usepackage{capt-of}

% ===== TikZ・図 =====
\usepackage{tikz}
\usetikzlibrary{positioning}
\usepackage{pgfplots}
\pgfplotsset{compat=1.18}

% ===== 化学 =====
\usepackage[version=4]{mhchem}
\usepackage{chemfig}

% ===== 回路 =====
\usepackage{circuitikz}

% ===== アルゴリズム =====
\usepackage{algorithm}
\usepackage{algpseudocode}

% ===== 数学補助 =====
% \numberwithin{equation}{section}
\renewcommand{\theequation}{\arabic{equation}}

\DeclareMathOperator{\sgn}{sgn}
\DeclareMathOperator{\tr}{tr}
\DeclareMathOperator{\diag}{diag}

\newcommand{\dif}{\mathrm{d}}
\newcommand{\ee}{\mathrm{e}}
\newcommand{\ii}{\mathrm{i}}

\newcommand{\pd}[2]{\frac{\partial{#1}}{\partial{#2}}}
\newcommand{\pdd}[2]{\frac{\partial^2{#1}}{\partial{#2}^2}}
\newcommand{\dd}[2]{\frac{\mathrm{d} #1}{\mathrm{d} #2}}

\newcommand{\vect}[1]{\bm{#1}}
\newcommand{\tensor}[1]{\bm{#1}}

% 画像がまだ配置されていない場合も，文書全体はコンパイルできるようにする．
% 指定した画像が存在すれば，通常どおりその画像を表示する．
\newcommand{\practiceimage}[2][\linewidth]{%
  \IfFileExists{#2}{%
    \includegraphics[width=#1]{#2}%
  }{%
    \fbox{%
      \parbox[c][28mm][c]{.9\dimexpr#1\relax}{%
        \centering\small
        画像ファイル\\未配置%
      }%
    }%
  }%
}

\lstset{
    language=C,
    basicstyle=\ttfamily\small,
    numbers=left,
    numberstyle=\tiny\color{gray},
    frame=single,
    rulecolor=\color{black},
    backgroundcolor=\color{gray!10},
    breaklines=true,
    showstringspaces=false,
    columns=fullflexible,
    keywordstyle=\color{blue},
    commentstyle=\color{green!50!black},
    stringstyle=\color{red},
    tabsize=4
}

\setlist[itemize]{topsep=4pt,itemsep=2pt,leftmargin=*}
\setlist[enumerate]{topsep=4pt,itemsep=2pt,leftmargin=*}
\setlength{\columnsep}{8mm}
\setlength{\emergencystretch}{1em}
\raggedbottom{}

\title{%
LaTeX演習用
% ここに自分で新しいテーマを付けてみる．テーマは自由である．
}
% ここを自分の名前に変更する．
\author{キムドホン}
% \date{\today}とすると，コンパイルした日付が入る．
\date{\today}

\begin{document}

\maketitle

\begin{abstract}
入退室管理システムが広く利用されている。
人物属性認識技術を組み合わせることで，帽子や眼鏡，バッグ，服装などの情報を取得できるため，人物照合の精度をさらに向上させることが可能となる。
このような複数の特徴情報を統合することで，従来の単一特徴に依存した認識方式よりも，高精度かつ安定した入退室管理システムの実現が期待される。
\end{abstract}
% 演習3：Abstractの挿入
% \begin{abstract} ... \end{abstract}を\maketitleの後へ追加する．

\section*{演習1：動作確認}
GitHubからこの「\texttt{LaTeX\_practice.tex}」をダウンロードし，
自分のPCでコンパイルしてみましょう．
正しくコンパイルできれば，PDFが生成されます．

\section*{演習2：テーマの変更}
タイトルと著者名を自分のものに変更し，\texttt{date}には\verb|\today|を入れて，
コンパイルしてみましょう．

\section*{演習3：Abstractの挿入}
Abstractを挿入することで，論文の概要を簡潔にまとめることができます．

\begin{verbatim}
\begin{abstract}
ここにAbstractを記述します．
\end{abstract}
\end{verbatim}

実際に，このコードを\verb|\maketitle|の後へ挿入してみましょう．

\section*{演習4：数式の挿入}
LaTeXでは，数式を美しく記述することができます．
数式の記述には，インライン数式とディスプレイ数式の2種類があります．

\subsection*{インライン数式}
円の面積は $S=\pi r^2$ で表される．

オイラーの公式は $e^{i\pi}+1=0$ である．

関数 $f(x)=x^2+2x+1$ を考える．

\subsection*{ディスプレイ数式}
\begin{equation}
\int_{0}^{1} x^2\,\dif x=\frac{1}{3}
\label{eq:integral}
\end{equation}

\begin{equation}
A=
\begin{pmatrix}
1 & 2\\
3 & 4
\end{pmatrix}
\label{eq:matrix}
\end{equation}

\begin{equation}
\lim_{n\to\infty}
{\left(1+\frac{1}{n}\right)}^n=e
\label{eq:euler}
\end{equation}

ディスプレイ数式は，\verb|\[ ... \]|で囲むと番号なしで表示できます．
番号を付けたい場合は，\verb|equation|環境を使用します．
実際に，インライン数式とディスプレイ数式をそれぞれ使って，
自分で数式を挿入してみましょう．

% ここに自分で数式を挿入する．

\subsection*{キムインライン数式}
球の表面積は $S=4\pi r^2$ で表される．

\subsection*{キムディスプレイ数式}
\begin{equation}
\oint_{S} \vec{E}\cdot \dif\vec{S}
=\frac{Q}{\varepsilon_0}
\label{eq:gauss}
\end{equation}


\section*{演習5：図の挿入}
LaTeXでは，図を挿入することもできます．

\begin{figure}[htbp]
    \centering
    \practiceimage[3cm]{1.jpg}
    \caption{LaTeXで挿入した図の例}\label{fig:custom-practice}
\end{figure}

画像の挿入は，画像ファイルの場所や拡張子，パッケージ設定などが原因で
エラーになることがあります．
エラーが出ても慌てず，エラーメッセージを読みながら一つずつ原因を確認しましょう．
では，自分で画像を挿入してみましょう．

% ここに画像を挿入する．
% 分からないときは，画像とともに生成AIへ
% 「画像を貼るLaTeXのコードを作成」などと質問してもよい．

\begin{figure}[htbp]
    \centering
    \includegraphics[width=3cm]{1.jpg}
    \caption{LaTeXで挿入した図の例}\label{fig:practice}
\end{figure}


\section*{演習6：表の挿入}
LaTeXでは，表を挿入することもできます．

\begin{table}[htbp]
    \centering
    \begin{tabular}{ccc}
        \toprule
        A & B & C \\
        \midrule
        1 & 2 & 3 \\
        4 & 5 & 6 \\
        \bottomrule
    \end{tabular}
    \caption{LaTeXで挿入した表の例}\label{tab:custom-practice}
\end{table}

では，自分で5行7列の表を挿入してみましょう．

% ここに表を挿入する．
% 分からないときは，表の画像とともに生成AIへ
% 「表を作るLaTeXのコードを作成」などと質問してもよい．

\begin{table}[htbp]
    \centering
    \begin{tabular}{cccc}
        \toprule
        A & B & C & D \\
        \midrule
        1 & 2 & 3 & 4 \\
        5 & 6 & 7 & 8 \\
        9 & 10 & 11 & 12 \\
        \bottomrule
    \end{tabular}
    \caption{LaTeXで挿入した表の例}\label{tab:practice}
\end{table}

\section*{演習7：箇条書き}
\begin{itemize}
    \item LaTeX
    \item Python
    \item GitHub
\end{itemize}

では，自分で箇条書きを作成してみましょう．

% ここに箇条書きを作成する．

\begin{itemize}
    \item キム
    \item ド
    \item ホン
\end{itemize}

\section*{演習8：ソースコードの挿入}
LaTeXでは，ソースコードもきれいに掲載できます．

\begin{lstlisting}
#include <stdio.h>

int main(void)
{
    for (int i = 0; i < 10; i++) {
        printf("%d\n", i);
    }

    return 0;
}
\end{lstlisting}

では，自分でソースコードを挿入してみましょう．

% ここにソースコードを挿入する．内容は自由である．

\begin{lstlisting}
#include <stdio.h>

int main(void)
{
    printf("Hello, kim do hoen\n");

    return 0;
}
\end{lstlisting}

\section*{演習9：URLの挿入}
URLを挿入することもできます．
OpenAIのホームページはこちらです．

\url{https://openai.com}

では，自分で任意のURLを挿入してみましょう．

% ここにURLを挿入する．

\url{https://www.naver.com}

\section*{演習10：参考文献}
卒業論文や学術論文では，他者の研究や論文を引用した場合，
必ず参考文献を明記する必要があります．

LaTeXでは，一般的にBibTeXを用いて参考文献を管理します．
参考文献を追加する手順は，大きく分けて次の3つです．

\subsection*{手順1：\texttt{.bib}ファイルを作成する}
まず，\texttt{references.bib}という名前のファイルを作成します．
このファイルには，引用する論文や書籍の情報を記述します．
例えば，次のように記述します．

\begin{lstlisting}[language={}]
@inproceedings{he2016resnet,
  author    = {Kaiming He and others},
  title     = {Deep Residual Learning for Image Recognition},
  booktitle = {CVPR},
  year      = {2016}
}
\end{lstlisting}

Google ScholarやIEEE Xploreでは，「BibTeX」を選択するだけで
この形式を取得できるため，自分で一から入力する必要はほとんどありません．

\subsection*{手順2：本文中で引用する}
引用したい箇所に，次のように記述します．

\begin{verbatim}
\cite{he2016resnet}
\end{verbatim}

例えば，

\begin{verbatim}
深層学習ではResNetが広く利用されている
\cite{he2016resnet}．
\end{verbatim}

と書くと，「深層学習ではResNetが広く利用されている[1]．」のように，
引用番号が自動で付与されます．

\subsection*{手順3：参考文献一覧を表示する}
文書の末尾へ，次の2行を追加します．

\begin{verbatim}
\bibliographystyle{ieeetr}
\bibliography{references}
\end{verbatim}

ここで，

\begin{itemize}
    \item \texttt{ieeetr}：参考文献の表示形式
    \item \texttt{references}：作成した\texttt{references.bib}の指定
\end{itemize}

を表しています．
実際に手順1から手順3までを行うと，本文では次のように表示されます．

「深層学習ではResNetが広く利用されている\cite{he2016resnet}．」

卒業論文では，参考文献の管理にBibTeXを用いることが一般的です．
今回は演習を省略しますが，研究を進める際には必ず利用する機能です．

\subsection*{発展：2段組で大きな図を挿入する}
卒業論文では，本文を2段組にし，図やグラフのみページ全体の幅で
表示することがあります．

その場合は，\verb|figure|の代わりに\verb|figure*|環境を使用します．
図がページ上部へ自動配置されることを確認するため，後ろにダミーテキストを入れています．

\begin{verbatim}
\begin{figure*}[t]
    \centering
    \includegraphics[width=0.8\textwidth]
        {1.jpg}
    \caption{実験結果}
    \label{fig:result}
\end{figure*}
\end{verbatim}

\begin{strip}
    \centering
    \begin{minipage}[t]{0.42\textwidth}
        \centering
        \includegraphics[width=0.72\linewidth]{1.jpg}
        \par\smallskip

        \small (a) 画像1
    \end{minipage}
    \hspace{0.06\textwidth}
    \begin{minipage}[t]{0.42\textwidth}
        \centering
        \includegraphics[width=0.72\linewidth]{1.jpg}
        \par\smallskip

        \small (b) 画像2
    \end{minipage}
    \captionof{figure}{画像を横並びにした例}\label{fig:images}
\end{strip}

図\ref{fig:images}では，同じ画像を2枚横に配置し，画像同士の大きさや
表示位置を比較しやすいレイアウトにしています．横並びの図を使用する場合は，
各画像の幅を揃え，画像間に適度な余白を設けると見やすくなります．

また，複数の画像を比較する図では，撮影条件や評価項目を統一することが重要です．
画像の下に短い説明を付けることで，読者は各画像の役割をすぐに把握できます．
本文では図を掲載するだけでなく，図から読み取れる特徴や比較結果についても
具体的に説明すると，資料全体の意図がより明確になります．

画像数が増える場合は，一つの画像を過度に小さくせず，必要に応じて図を分けます．
ページ内の余白，本文との間隔，キャプションの位置を揃えることで，
情報量が多い場合でも読みやすいレイアウトを維持できます．

% 参考文献一覧は，文書の最後に一度だけ出力する．
\bibliographystyle{ieeetr}
\bibliography{references}

\end{document}
